% strat_t5_costs — turnover and net-of-transaction-cost performance.
% Numbers transcribed from the output of R/strategy_ext.R.
\small
\begin{tabular}{l c c c c c}
  \toprule
  & Turnover & \multicolumn{4}{c}{One-way cost (bp)} \\
  \cmidrule(lr){3-6}
  Strategy & (\%/month) & $0$ & $10$ & $25$ & $50$ \\
  \midrule
  \panel{6}{Panel A: Sharpe ratio} \\
  $\GCF$-timed           & $9.1$ & $0.41$ & $0.39$ & $0.37$ & $0.34$ \\
  Recursive-mean timing  & $1.3$ & $0.37$ & $0.37$ & $0.36$ & $0.36$ \\
  Buy-and-hold           & --    & $0.34$ & $0.34$ & $0.34$ & $0.34$ \\
  \midrule
  \panel{6}{Panel B: CER, $\gamma=5$ (\%)} \\
  $\GCF$-timed           & $9.1$ & $1.66$ & $1.55$ & $1.39$ & $1.11$ \\
  Recursive-mean timing  & $1.3$ & $1.35$ & $1.33$ & $1.31$ & $1.27$ \\
  Buy-and-hold           & --    & $1.16$ & $1.16$ & $1.16$ & $1.16$ \\
  \bottomrule
\end{tabular}
\tabnotes{This table reports the strategies of \Cref{tab:strat-perf} net of
  proportional transaction costs. Turnover is the average one-way monthly
  turnover $|w_{t}-w_{t-1}|$ of the equal-average-exposure weight path, in
  percent of the portfolio notional; the annualised cost drag $c\cdot 12\cdot
  |w_{t}-w_{t-1}|$ is subtracted from each twelve-month return for a one-way
  cost of $c$ basis points on the notional traded. Buy-and-hold is assumed
  costless after the initial purchase. The first month of the evaluation
  window is dropped (turnover undefined), so the zero-cost column differs
  marginally from \Cref{tab:strat-perf}. For reference, quoted half-spreads on
  ten-year G10 government bond futures are of the order of one basis point,
  and even for cash bonds one-way costs above $25$ basis points are
  conservative.}
